\section{Marked completion and the integer Euler obstruction}
\label{sec:euler}

\subsection{Relative complexes and a persistent marked point}

Fix a classical one-component planar diagram $D$ with $N$ crossings and a
validated permutation of its crossings. A proper prefix leaves a suffix with
$n\ge1$ crossings. The scanning construction represents the prefix by a complex
in the dotted cobordism category. Each object is a boundary matching with its
homological and quantum shifts; closed circles have been removed through the
standard delooping operations. Differentials, cancellations, and their homotopies
are the actual maps supplied by this construction \cite{barnatan2005,barnatan2007}.

Choose a point on an edge incident to the final unprocessed crossing. It remains
in the suffix at every proper-prefix observation. On completion, the marked
circle carries the action of $A=\F_2[x]/(x^2)$. The marked quotient is applied
after the actual suffix has been attached.

\begin{lemma}[Why completion and reduction preserve the contractions]
\label{lem:marked}
Suppose the prefix complex is decomposed and simplified by chain maps and
chain homotopies that commute with the marked action after completion. Applying
the marked quotient preserves every recorded chain-homotopy equivalence.
\end{lemma}
\begin{proof}
If maps $f,g$ and a homotopy $h$ satisfy $fg-1=dh+hd$, all maps are
$A$-linear. Their induced maps on the quotient satisfy the same equality,
because passage to the quotient respects composition and addition. The other
homotopy identity is treated identically. This argument does not require the
quotient functor to be flat or exact on arbitrary short exact sequences.
\end{proof}

Let $\mathcal C_a$ be a whole connected component of the differential graph after such
simplification. The absence of differentials between these groups gives a
direct sum of complexes. Completing with the actual suffix and applying the
marked quotient gives
\begin{equation}
 \Khred(K;\F_2)\cong
 \bigoplus_a H\!\left(\widetilde{\operatorname{cl}}(\mathcal C_a)\right)^{\oplus m_a},
 \label{eq:completed-sum}
\end{equation}
up to recorded shifts, whose effect on the norms below is harmless. Shared
copies can have different global shifts; each has the same norm.

\subsection{Four integer coordinates}

For a finite bigraded complex $X$, define its residue-four Euler vector by
\begin{equation}
 \sigma_r(X)=\sum_{h,\ j\equiv r\pmod4}(-1)^h\dim_{\F_2}X^{h,j},
 \qquad 0\le r<4.
 \label{eq:sigma}
\end{equation}
Differentials preserve the quantum grading, so taking Euler characteristics
after homology leaves this vector unchanged. A homological shift negates all
four coordinates when the shift is odd. A quantum shift cyclically permutes
them. Both preserve the $\ell^1$ norm.

\begin{proposition}[Existing marked-component bound]
\label{prop:B4}
With the provenance of \cref{lem:marked}, put
$x_a=\sigma(\widetilde{\operatorname{cl}}(\mathcal C_a))$. Then
\[
 B_4=\sum_a m_a\norm{x_a}
 \le\dim_{\F_2}\Khred(K;\F_2).
\]
In particular, $B_4>1$ certifies that $K$ is nontrivial.
\end{proposition}
\begin{proof}
For each residue, let $e_r,o_r$ be the total dimensions of the even and odd
homology groups. Then $|e_r-o_r|\le e_r+o_r$. Summing over residues and
using \cref{eq:completed-sum} proves the inequality. Reduced Khovanov rank one
detects the unknot; the maintained $\F_2$ criterion follows from the standard
unknot-detection theorem and change-of-coefficients implications
\cite{kronheim}. The rank-one implication is the existing recognition premise,
not a new assertion about the auxiliary prime fields used below.
\end{proof}

\paragraph{The order of operations is part of the theorem.}
For each whole component, sum all signed, shifted object vectors first. Then
take a norm, and finally sum with positive multiplicities. A pair of object
vectors $u$ and $-u$ may belong to one contractible complex. Their true Euler
contribution is zero, while $\norm{u}+\norm{-u}$ is positive. Neither modular
arithmetic nor early stopping permits taking a norm midway through this sum.
Stopping after a \emph{completed} component is safe because all remaining
component norms are nonnegative.

\subsection{Recovering the vector from two scalar evaluations}

We use the repository's raw reduced polynomial convention
\begin{equation}
 J_D(q)=\sum_{s\in\{0,1\}^n}(-q)^{|s|}
                 (q+q^{-1})^{c(s)-1}.
 \label{eq:raw-J}
\end{equation}
Here $c(s)$ counts the circles of the smoothing. The completion contains no
explicit crossing-free circles outside this diagram; previously delooped
circles are represented by the prefix's shifted objects and multiplicities.

Changing a smoothing changes the circle count by one in parity. Hence all
quantum degrees in \cref{eq:raw-J} have parity
$\epsilon=(c(0)-1)\bmod2$. If
\[
 a=J_D(1),\qquad d=i^{-\epsilon}J_D(i)\in\Z,
\]
the only potentially nonzero coordinates are
\begin{equation}
 \sigma_\epsilon=\frac{a+d}{2},\qquad
 \sigma_{\epsilon+2}=\frac{a-d}{2}.
 \label{eq:two-values}
\end{equation}
Indeed, evaluation at one adds the two residue sums, while evaluation at $i$
subtracts them after removing the common factor $i^\epsilon$. In particular,
\begin{equation}
                     \norm{\sigma(D)}=\max(|a|,|d|).
 \label{eq:maxnorm}
\end{equation}
This identity will give small and explicit coefficient bounds.

\subsection{The exact evaluation at one}

Let $c$ be the number of link components and $c_0=c(0)$ the all-zero circle
count. In the raw convention,
\begin{equation}
                  a=(-1)^{n+c+c_0}2^{c-1}.
 \label{eq:euler-sign}
\end{equation}
One way to see the sign is to choose an orientation and write $h$ for the
number of one-smoothings in its Seifert state and $s$ for its Seifert circles.
Circle parity gives $s\equiv c_0+h\pmod2$. The orientable Seifert surface
has $s-n\equiv c\pmod2$. Thus $h\equiv n+c+c_0\pmod2$.
The normalized reduced evaluation at one is $2^{c-1}$, and the raw shift
supplies the sign above.

The existing implementation computes the same integer by independently
orienting the completed link. A reusable boundary connectivity summary can
instead determine $c$ and $c_0$ without revisiting the suffix. Both versions
remain available as cross-checks. The integer $2^{c-1}$ has $O(n)$ bits; a
boundary-sized combinatorial traversal is not automatically a boundary-only
bit-cost bound for constructing this integer.

\subsection{The signed determinant and its phase}

For a connected completed projection, checkerboard color its regions. At a
crossing $e$, let $b_e$ be the smoothing bit that does not join its black
sectors, and put $w_e=(-1)^{b_e}$. Set $B=\sum_e b_e$. If $G$ is the signed
black Tait graph with $v$ vertices, then
\begin{equation}
 J_D(i)=(-i)^{B+v-1}\tau(G),\qquad
 \tau(G)=\det L_G^{(0)},
 \label{eq:tait}
\end{equation}
where $L_G^{(0)}$ is a grounded Laplacian cofactor.

At $q=i$, all states except one-circle states vanish. Those states correspond
to spanning trees of the Tait graph. Their phases give \cref{eq:tait}; the
matrix-tree expansion is polynomial in the edge weights, so the signs do not
require a positive Laplacian. Loops contribute to $B$, although their incidence
vectors vanish in the Laplacian. Removing them from the phase would invalidate
the observation. If the projection is disconnected, every smoothing has more
than one circle and $J_D(i)=0$.

For connected accepted geometry we also have
\[
                       B+v-1\equiv\epsilon\pmod2.
\]
Choose any unsigned spanning tree. Its one-circle state has smoothing-weight
parity $B+v-1$, while circle parity gives $|s|\equiv c_0-1\pmod2$.
This proves the identity even when the signed tree sum cancels.
It is a combinatorial parity condition, not merely a check that the computed
Gaussian integer lies on the correct axis. The implementation checks it even
when a modular determinant vanishes, so a zero residue cannot conceal a phase
inconsistency.

\subsection{Relative shifts and multiplicity saturation}

The maintained grading recovery assigns a relative quantum shift $t_o$ to
every object in a component, checking every nonzero differential and cycle.
For a map of dot degree $d$ between matchings with $m$ arcs and $c$ glued
circles, its equation is $t_{\rm target}-t_{\rm source}=m-c+2d$.
Object $o$ contributes $(-1)^{h_o}\operatorname{rot}_{t_o}\sigma(D_o)$.
Only the shift modulo four is needed after consistency has been verified.

For threshold $T$, replacing a positive multiplicity $m$ by
$\min(m,T+1)$ preserves whether $m\norm{x}>T$, because the norm is an integer.
The same statement extends to sums of nonnegative terms. The maintained
scanner uses cap three; this remains sound for our threshold one. It does not
permit saturating signed Euler coefficients before aggregation.

Finally, a proper-prefix restriction remains mandatory. At the closed terminal
stage, separate unmarked scalar generators are not automatically free marked
$A$-summands. The modular observer refuses that stage. An \code{UNKNOT} verdict
comes only from the completed rank-two unreduced calculation.
